\subsection*{Manuscript preparation}
Large language models (OpenAI ChatGPT, Google Gemini and Anthropic Claude) assisted with manuscript drafting and language editing and with the formulation and checking of derivations, analysis scripts and figure code. The supplied plotting scripts draw the system diagrams from the stated model and the experimental panels from the accompanying data. The authors take responsibility for the final content.

\section*{Data availability}
The Source Data archive accompanies this Article and is also available to reviewers with access to \url{https://github.com/pooya-hejazi/uplink-scheduling}. It contains the numerical values underlying Figs.~1--6, Supplementary Figs.~1--8 and Supplementary Tables~1--16. The data supplied for reproduction include the simulated channels, frozen schedules, power bounds, matching witnesses and outcomes for individual channel realizations. The O1\_60 and ASU Campus 1 ray-tracing scenarios are available from the \href{https://www.deepmimo.net/}{DeepMIMO database}~\cite{deepmimo}. The supplied manifests identify the base station, spatial tiles and users used for each scene. The six-user distributed recordings are available from the \href{https://renew-wireless.org/dataset-distributed.html}{Rice RENEW distributed collection} and its IEEE DataPort record~\cite{zafari}. The UHF recordings are available through the \href{https://renew-wireless.org/dataset-argos.html}{Rice Argos Channel Survey}~\cite{shepard}. Access to the original radio samples requires registration with Rice and acceptance of its research-use terms. The supplied files give capture identifiers, selected byte ranges, decoding conventions and derived measurements.
Pooya Hejazi maintains the repository and is the contact for questions about reproducing the results.


\section*{Code availability}
Reviewers can obtain the code in the accompanying Supplementary Software~1 or request access to the private repository at \url{https://github.com/pooya-hejazi/uplink-scheduling}. Pooya Hejazi is the contact for repository access and reproduction questions. The supplied software includes the scheduler, the certificate checker, the C++ matching interface, the separate CPU verification re-implementation, the source-data analysis and the instructions needed to reproduce the figures. Matching uses the LEMON graph library~\cite{lemon}. The simulation and measurement processing use \href{https://github.com/NVlabs/sionna}{Sionna}, \href{https://github.com/DeepMIMO/DeepMIMO}{DeepMIMO} and decoding conventions from \href{https://github.com/renew-wireless/RENEWLab}{RENEWLab}.

\begin{thebibliography}{51}\small
\bibitem{marzetta} Marzetta, T. L. \href{https://doi.org/10.1109/TWC.2010.092810.091092}{Noncooperative cellular wireless with unlimited numbers of base station antennas}. \emph{IEEE Trans. Wirel. Commun.} \textbf{9}, 3590–3600 (2010).
\bibitem{ngo} Ngo, H. Q., Ashikhmin, A., Yang, H., Larsson, E. G. \& Marzetta, T. L. \href{https://doi.org/10.1109/TWC.2017.2655515}{Cell-free massive MIMO versus small cells}. \emph{IEEE Trans. Wirel. Commun.} \textbf{16}, 1834–1850 (2017).
\bibitem{bjornson} Björnson, E. \& Sanguinetti, L. \href{https://doi.org/10.1109/TWC.2019.2941478}{Making cell-free massive MIMO competitive with MMSE processing and centralized implementation}. \emph{IEEE Trans. Wirel. Commun.} \textbf{19}, 77–90 (2020).
\bibitem{demir} Demir, Ö. T., Björnson, E. \& Sanguinetti, L. \href{https://doi.org/10.1561/2000000109}{Foundations of user-centric cell-free massive MIMO}. \emph{Found. Trends Signal Process.} \textbf{14}, 162–472 (2021).
\bibitem{buzzi} Buzzi, S., D’Andrea, C., Zappone, A. \& D’Elia, C. \href{https://doi.org/10.1109/TWC.2019.2952117}{User-centric 5G cellular networks: resource allocation and comparison with the cell-free massive MIMO approach}. \emph{IEEE Trans. Wirel. Commun.} \textbf{19}, 1250–1264 (2020).
\bibitem{caire} Caire, G. \href{https://doi.org/10.1109/TWC.2018.2808522}{On the ergodic rate lower bounds with applications to massive MIMO}. \emph{IEEE Trans. Wirel. Commun.} \textbf{17}, 3258–3268 (2018).
\bibitem{wang} Wang, Z., Zhang, J., Ngo, H. Q., Ai, B. \& Debbah, M. \href{https://doi.org/10.1109/TCOMM.2023.3235919}{Uplink precoding design for cell-free massive MIMO with iteratively weighted MMSE}. \emph{IEEE Trans. Commun.} \textbf{71}, 1646–1664 (2023).
\bibitem{haustein} Haustein, T. \href{https://edocs.tib.eu/files/e01dd01/522593321.pdf}{Real time signal processing for multi-antenna systems and experimental verification on a reconfigurable hardware test-bed}. \emph{PhD thesis, Technische Universität Berlin} (2006).
\bibitem{viterbo} Viterbo, E. \& Hottinen, A. \href{https://doi.org/10.1109/ISSSTA.2008.50}{Optimal user pairing for multiuser MIMO}. \emph{Proc. 10th IEEE Int. Symp. Spread Spectrum Tech. Appl. (ISSSTA)}, 242–246 (2008).
\bibitem{fan} Fan, J., Li, G. Y., Yin, Q., Peng, B. \& Zhu, X. \href{https://doi.org/10.1109/TWC.2012.070912.111305}{Joint user pairing and resource allocation for LTE uplink transmission}. \emph{IEEE Trans. Wirel. Commun.} \textbf{11}, 2838–2847 (2012).
\bibitem{ruder2013} Ruder, M. A., Ding, D., Dang, U. L., Vasilakos, A. V. \& Gerstacker, W. H. \href{https://doi.org/10.1016/j.phycom.2013.02.004}{Joint user grouping and frequency allocation for multiuser SC-FDMA transmission}. \emph{Phys. Commun.} \textbf{8}, 91–103 (2013).
\bibitem{ruder2014} Ruder, M. A., Moldovan, A. \& Gerstacker, W. H. \href{https://doi.org/10.1016/j.phycom.2014.10.001}{SC-FDMA user pairing and frequency allocation with imperfect channel state information}. \emph{Phys. Commun.} \textbf{13}, 88–98 (2014).
\bibitem{prasad} Prasad, N., Zhang, H., Zhu, H. \& Rangarajan, S. \href{https://doi.org/10.1109/TWC.2013.072513.120734}{Multi-user MIMO scheduling in the fourth generation cellular uplink}. \emph{IEEE Trans. Wirel. Commun.} \textbf{12}, 4272–4285 (2013).
\bibitem{viswanath} Viswanath, P., Tse, D. N. C. \& Anantharam, V. \href{https://doi.org/10.1109/18.904525}{Asymptotically optimal water-filling in vector multiple-access channels}. \emph{IEEE Trans. Inf. Theory} \textbf{47}, 241–267 (2001).
\bibitem{shi} Shi, S., Schubert, M. \& Boche, H. \href{https://doi.org/10.1109/TSP.2008.924796}{Rate optimization for multiuser MIMO systems with linear processing}. \emph{IEEE Trans. Signal Process.} \textbf{56}, 4020–4030 (2008).
\bibitem{shenfpone} Shen, K. \& Yu, W. \href{https://doi.org/10.1109/TSP.2018.2812733}{Fractional programming for communication systems—Part I: Power control and beamforming}. \emph{IEEE Trans. Signal Process.} \textbf{66}, 2616–2630 (2018).
\bibitem{shenfptwo} Shen, K. \& Yu, W. \href{https://doi.org/10.1109/TSP.2018.2812748}{Fractional programming for communication systems—Part II: Uplink scheduling via matching}. \emph{IEEE Trans. Signal Process.} \textbf{66}, 2631–2644 (2018).
\bibitem{sun} Sun, H. et al. \href{https://doi.org/10.1109/TSP.2018.2866382}{Learning to optimize: training deep neural networks for interference management}. \emph{IEEE Trans. Signal Process.} \textbf{66}, 5438–5453 (2018).
\bibitem{shen} Shen, Y., Shi, Y., Zhang, J. \& Letaief, K. B. \href{https://doi.org/10.1109/JSAC.2020.3036965}{Graph neural networks for scalable radio resource management: architecture design and theoretical analysis}. \emph{IEEE J. Sel. Areas Commun.} \textbf{39}, 101–115 (2021).
\bibitem{hejazicluster} Hejazi, P. \& Sabaei, M. \href{https://doi.org/10.23919/JCN.2025.000085}{Optimizing spectral efficiency through user clustering in user-centric cell-free massive MIMO environment for 6G wireless network}. \emph{J. Commun. Netw.} \textbf{28}, 30–39 (2026).
\bibitem{ahmadi2026} Ahmadi, N. \& Akbarizadeh, G. \href{https://doi.org/10.1038/s41598-026-38685-3}{Machine learning based power control in cellular and cell-free massive MIMO systems}. \emph{Sci. Rep.} \textbf{16}, 8129 (2026).
\bibitem{hamakawa2026} Hamakawa, Y., Kashimata, T., Yamasaki, M. \& Tatsumura, K. \href{https://doi.org/10.1038/s41467-026-73725-6}{Machine learning-assisted high-speed combinatorial optimization with Ising machines for dynamically changing problems}. \emph{Nat. Commun.} \textbf{17}, 4877 (2026).
\bibitem{kashimata2026} Kashimata, T., Hamakawa, Y., Yamasaki, M. \& Tatsumura, K. \href{https://doi.org/10.1038/s41467-026-76069-3}{Real-time black-box optimization for dynamic discrete environments using embedded Ising machines}. \emph{Nat. Commun.} \textbf{17}, 8085 (2026).
\bibitem{dinitz} Dinitz, M., Im, S., Lavastida, T., Moseley, B. \& Vassilvitskii, S. \href{https://proceedings.neurips.cc/paper/2021/hash/5616060fb8ae85d93f334e7267307664-Abstract.html}{Faster matchings via learned duals}. \emph{Adv. Neural Inf. Process. Syst.} \textbf{34}, 10393–10406 (2021).
\bibitem{bessa} Bessa, S., Dabert, D., Bourgeat, M., Rousseau, L.-M. \& Cappart, Q. \href{https://doi.org/10.1609/aaai.v39i11.33208}{Learning valid dual bounds in constraint programming: boosted Lagrangian decomposition with self-supervised learning}. \emph{Proc. AAAI Conf. Artif. Intell.} \textbf{39}, 11113–11121 (2025).
\bibitem{edmonds} Edmonds, J. \href{https://doi.org/10.6028/jres.069B.013}{Maximum matching and a polyhedron with 0,1-vertices}. \emph{J. Res. Natl Bur. Stand. B} \textbf{69B}, 125–130 (1965).
\bibitem{kuhn} Kuhn, H. W. \href{https://doi.org/10.1002/nav.3800020109}{The Hungarian method for the assignment problem}. \emph{Nav. Res. Logist. Q.} \textbf{2}, 83–97 (1955).
\bibitem{tr} 3GPP. \href{https://www.etsi.org/deliver/etsi_tr/138900_138999/138901/19.02.00_60/tr_138901v190200p.pdf}{Study on channel model for frequencies from 0.5 to 100 GHz}. \emph{3GPP TR 38.901}, V19.2.0, Release 19 (2026).
\bibitem{sionna} Hoydis, J. et al. \href{https://doi.org/10.48550/arXiv.2203.11854}{Sionna: an open-source library for next-generation physical layer research}. \emph{arXiv}, 2203.11854 (2022).
\bibitem{deepmimo} Alkhateeb, A. \href{https://arxiv.org/abs/1902.06435}{DeepMIMO: a generic deep learning dataset for millimeter wave and massive MIMO applications}. \emph{Proc. Inf. Theory Appl. Workshop (ITA)}, 1–8 (2019).
\bibitem{zafari} Zafari, M. \href{https://doi.org/10.21227/74ez-b275}{Distributed Multi-User MIMO Datasets}. \emph{IEEE DataPort} (2024).
\bibitem{renew} Doost-Mohammady, R. et al. \href{https://doi.org/10.1109/ACSSC.2018.8645391}{RENEW: programmable and observable massive MIMO networks}. \emph{Proc. 52nd Asilomar Conf. Signals Syst. Comput.}, 1654–1658 (2018).
\bibitem{barnhart} Barnhart, C., Johnson, E. L., Nemhauser, G. L., Savelsbergh, M. W. P. \& Vance, P. H. \href{https://doi.org/10.1287/opre.46.3.316}{Branch-and-price: column generation for solving huge integer programs}. \emph{Oper. Res.} \textbf{46}, 316–329 (1998).
\bibitem{matthiesen} Matthiesen, B., Hellings, C., Jorswieck, E. A. \& Utschick, W. \href{https://doi.org/10.1109/TSP.2020.2983284}{Mixed monotonic programming for fast global optimization}. \emph{IEEE Trans. Signal Process.} \textbf{68}, 2529–2544 (2020).
\bibitem{yului} Yu, W. \& Lui, R. \href{https://doi.org/10.1109/TCOMM.2006.877962}{Dual methods for nonconvex spectrum optimization of multicarrier systems}. \emph{IEEE Trans. Commun.} \textbf{54}, 1310–1322 (2006).
\bibitem{zhaopottie} Zhao, Y. \& Pottie, G. J. \href{https://doi.org/10.1109/TIT.2013.2255731}{Optimal spectrum management in multiuser interference channels}. \emph{IEEE Trans. Inf. Theory} \textbf{59}, 4961–4976 (2013).
\bibitem{dyer} Dyer, M. E. \& Frieze, A. M. \href{https://doi.org/10.1016/0196-6774(86)90002-7}{Planar 3DM is NP-complete}. \emph{J. Algorithms} \textbf{7}, 174–184 (1986).
\bibitem{razaviyayn} Razaviyayn, M., Sanjabi, M. \& Luo, Z.-Q. \href{https://doi.org/10.1109/TIT.2012.2184909}{Linear transceiver design for interference alignment: complexity and computation}. \emph{IEEE Trans. Inf. Theory} \textbf{58}, 2896–2910 (2012).
\bibitem{salaun} Salaün, L., Chen, C. S. \& Coupechoux, M. \href{https://doi.org/10.1109/ICC.2018.8422362}{Optimal joint subcarrier and power allocation in NOMA is strongly NP-hard}. \emph{Proc. IEEE Int. Conf. Commun. (ICC)}, 1–7 (2018).
\bibitem{hejazipilot} Hejazi, P., Sabaei, M. \& Hajmahmoudi, A. \href{https://doi.org/10.23919/JCN.2025.000112}{Pilot-efficient channel estimation for 6G MU-mMIMO using LS initialization and GAN-augmented CNN refinement}. \emph{J. Commun. Netw.} \textbf{28}, 401–408 (2026).
\bibitem{lei} Lei, J., G’Sell, M., Rinaldo, A., Tibshirani, R. J. \& Wasserman, L. \href{https://doi.org/10.1080/01621459.2017.1307116}{Distribution-free predictive inference for regression}. \emph{J. Am. Stat. Assoc.} \textbf{113}, 1094–1111 (2018).
\bibitem{johnstone} Johnstone, C. \& Cox, B. \href{https://proceedings.mlr.press/v152/johnstone21a.html}{Conformal uncertainty sets for robust optimization}. \emph{Proc. Mach. Learn. Res.} \textbf{152}, 72–90 (2021).
\bibitem{kludze2026} Kludze, A., Bodet, D., Jornet, J. M. \& Ghasempour, Y. \href{https://doi.org/10.1038/s41467-026-71793-2}{Distributed wavefront shaping in radiative near-field sub-terahertz wireless networks}. \emph{Nat. Commun.} \textbf{17}, 5595 (2026).
\bibitem{shenoam2026} Shen, R. \& Ghasempour, Y. \href{https://doi.org/10.1038/s41467-026-69493-y}{Joint communication and sensing with structured beams carrying orbital angular momentum}. \emph{Nat. Commun.} \textbf{17}, 2832 (2026).
\bibitem{rogalin} Rogalin, R. et al. \href{https://doi.org/10.1109/TWC.2014.030314.130474}{Scalable synchronization and reciprocity calibration for distributed multiuser MIMO}. \emph{IEEE Trans. Wirel. Commun.} \textbf{13}, 1815–1831 (2014).
\bibitem{vieira} Vieira, J. et al. \href{https://doi.org/10.1109/TWC.2017.2674659}{Reciprocity calibration for massive MIMO: proposal, modeling, and validation}. \emph{IEEE Trans. Wirel. Commun.} \textbf{16}, 3042–3056 (2017).
\bibitem{ganesan} Ganesan, U. K., Sarvendranath, R. \& Larsson, E. G. \href{https://doi.org/10.1109/TWC.2023.3335089}{BeamSync: over-the-air synchronization for distributed massive MIMO systems}. \emph{IEEE Trans. Wirel. Commun.} \textbf{23}, 6824–6837 (2024).
\bibitem{zhangsync2026} Zhang, C. et al. \href{https://doi.org/10.1038/s41467-026-76739-2}{Sub-picosecond synchronized multi-faceted fronthaul with one-sample-per-symbol coherent detection}. \emph{Nat. Commun.} \textbf{17}, 9747 (2026).
\bibitem{ni2026} Ni, Z. et al. \href{https://doi.org/10.1038/s41467-026-74485-z}{Fully-connected microwave photonic multi-beamformer with fast beam-steering for broadband wireless communication}. \emph{Nat. Commun.} \textbf{17}, 7743 (2026).
\bibitem{lemon} Dezső, B., Jüttner, A. \& Kovács, P. \href{https://doi.org/10.1016/j.entcs.2011.06.003}{LEMON – an open source C++ graph template library}. \emph{Electron. Notes Theor. Comput. Sci.} \textbf{264}, 23–45 (2011).
\bibitem{shepard} Shepard, C., Ding, J., Guerra, R. E. \& Zhong, L. \href{https://doi.org/10.1109/ACSSC.2016.7869082}{Understanding real many-antenna MU-MIMO channels}. \emph{Proc. 50th Asilomar Conf. Signals Syst. Comput.}, 461–467 (2016).
\end{thebibliography}

\section*{Acknowledgements}
This study received no funding. We acknowledge use of the Rice RENEW distributed recordings, the Rice Argos Channel Survey and the DeepMIMO ray-tracing scenarios.
\section*{Author contributions}
P.H. conceived the study, formulated the mathematical bounds, developed the scheduling and verification software, analysed the simulation and measurement data, and wrote the manuscript. M.S. supervised the research, contributed to the problem formulation and analytical design, and reviewed and edited the manuscript.

\section*{Competing interests}
The authors declare no competing interests.

